% Fiche Python 3 (version avancée) — recto-verso A4 paysage
% Source : pas de cours ; complément de python_cheat.tex (bases, fichiers, numpy...) :
%   ici itérateurs, générateurs, décorateurs, objets, context managers, typage,
%   concurrence (threads, processus, asyncio), tests, débogage et pièges avancés.
% Compiler : make (ici ou à la racine), ou xelatex python_avance.tex — une seule passe suffit
% Apparence : ../librairie (fiche.sty + theme-python.tex, partagé avec python_cheat.tex).
%   colonne[titre][couleur] (3 par page, séparées par \colbreak) · bloc{titre}{cat}
%   catégories : syn typ ctl ex io lib (couleurs et icônes dans le thème)
%   code · \V|...| code en ligne · «...» dans le code = à remplir
%   \sub sous-titre · \note remarque · \piege piège · \badge{TXT} · schema
%   optab étiquette | texte · tab3 étiquette | code | remarque (largeurs mesurées)
% Le code est en police à chasse fixe : les commentaires sont alignés à la main
% (attention : les « » ne sont pas affichés, ils ne comptent pas dans l'alignement).
% Après compilation : grep "\[fiche\]" python_avance.log (place libre, collisions).
\documentclass{article}
\usepackage[python]{../librairie/fiche}   % apparence : ../librairie (fiche.sty + theme-python.tex)
% --- Langage Python 3 pour listings ---
% Classes de mots-clés : [1] rose · [2] cyan · [3] vert · [4] couleur du thème
% (styles dans ../librairie/fiche.sty). Chargé par python_avance.tex.
% [1] mots-clés · [2] types · [3] fonctions / modules / décorateurs intégrés
\lstdefinelanguage{CODE}{
  sensitive=true, alsoletter={@},
  morekeywords=[1]{False,None,True,and,as,assert,async,await,break,class,continue,def,
    del,elif,else,except,finally,for,from,global,if,import,in,is,lambda,nonlocal,not,
    or,pass,raise,return,try,while,with,yield,match,case,type},
  morekeywords=[2]{int,float,complex,bool,str,bytes,list,tuple,dict,set,frozenset,
    range,object,Exception,ValueError,KeyError,TypeError,StopIteration,IndexError,
    OSError,FileNotFoundError,
    Iterator,Iterable,Generator,Callable,Optional,Protocol,TypedDict,Literal,Final,
    TypeVar,Any,Self,NamedTuple,Enum,IntFlag,ABC,Path},
  morekeywords=[3]{print,len,enumerate,zip,map,filter,sorted,sum,min,max,iter,next,
    isinstance,super,getattr,setattr,repr,hash,format,self,cls,
    @dataclass,@property,@staticmethod,@classmethod,@contextmanager,
    @cache,@lru_cache,@abstractmethod,@total_ordering,@singledispatch,
    @cached_property,@timer,@retry,@c,
    __init__,__repr__,__str__,__eq__,__hash__,__lt__,__len__,__getitem__,
    __iter__,__next__,__contains__,__call__,__enter__,__exit__,__add__,
    __post_init__,__name__,__slots__,__mro__,
    asyncio,itertools,functools,contextlib,dataclasses,typing,pytest,logging,
    subprocess,time,os,sys,pdb,traceback},
  morecomment=[l]{\#}, morestring=[b]", morestring=[b]',
}
   % coloration syntaxique : \lstdefinelanguage{CODE}

\begin{document}
%%%%%%%%%%%%%%%%%%%%%%%%%%%%%%%%%%%%%%%%%%%%%%%%%%%%%%%%%%%%%%%%%%%%%%%%%%%%%%
%                         RECTO : LE LANGAGE AVANCÉ
%%%%%%%%%%%%%%%%%%%%%%%%%%%%%%%%%%%%%%%%%%%%%%%%%%%%%%%%%%%%%%%%%%%%%%%%%%%%%%
\entete{Python avancé}{recto · itérateurs, fonctions, objets}
\noindent
% ======================= colonne 1 : match, itérer =======================
\begin{colonne}[match · itérer][col-a]
\begin{bloc}{match : motifs avancés}{ctl}
\begin{code}
match msg:
    case {"op": "write", "adr": int(a), "data": d}:
        ecrire(a, d)                  # motif de dict
    case Point(x=0, y=y):             # motif de classe
        ...
    case [tete, *reste] if tete > 0:  # garde
        ...
    case int() | float() as v:        # alternative + nom
        ...
    case _:
        raise ValueError(msg)
\end{code}
\end{bloc}

\begin{bloc}{Itérateurs}{ctl}
\begin{code}
it = iter([1, 2, 3])  # appelle __iter__
next(it), next(it)    # 1, 2 (appelle __next__)
next(it, None)        # défaut, pas d'exception
class Compte:         # itérateur à la main
    def __init__(self, n): self.i, self.n = 0, n
    def __iter__(self): return self
    def __next__(self):
        if self.i >= self.n: raise StopIteration
        self.i += 1
        return self.i
\end{code}
\note{Itérateur : \V|__iter__| + \V|__next__| ; ne se parcourt qu'une fois.}
\end{bloc}

\begin{bloc}{Générateurs}{ctl}
\begin{code}
def lire_mots(chemin):       # génératrice
    with open(chemin) as fh:
        for ligne in fh:
            yield from ligne.split()  # délègue
mots = lire_mots("log.txt")  # rien n'est lu ici
for m in mots: ...           # lu à la demande
def fib():
    a, b = 0, 1
    while True:              # infini : sans souci
        yield a
        a, b = b, a + b
carres = (x * x for x in range(10**9))  # paresseux
\end{code}
\note{\V|yield| suspend la fonction ; l'appel suivant reprend juste après. Mémoire constante.}
\end{bloc}

\begin{bloc}{itertools}{lib}
\begin{optab}[27mm]
\V|islice(it, 5)|      & 5 premiers d'un itérateur \\
\V|count(0, 2)|        & 0, 2, 4... à l'infini \\
\V|chain(l1, l2)|      & enchaîner \\
\V|product(r, r)|      & produit cartésien (boucles imbriquées) \\
\V|combinations(l, 2)| & paires, sans ordre ni répétition \\
\V|permutations(l)|    & tous les ordres \\
\V|groupby(l, key)|    & groupes consécutifs (trier avant) \\
\V|accumulate(l)|      & sommes cumulées \\
\V|pairwise(l)|        & (a, b), (b, c)... \badge{3.10} \\
\V|batched(l, 4)|      & paquets de 4 \badge{3.12} \\
\V|zip_longest(a, b)|  & zip jusqu'au plus long
\end{optab}
\end{bloc}
\end{colonne}\colbreak
% ======================= colonne 2 : fonctions =======================
\begin{colonne}[fonctions][col-b]
\begin{bloc}{Fermetures (closures)}{syn}
\begin{code}
def compteur(debut=0):
    n = debut
    def inc(pas=1):
        nonlocal n        # modifier n du dessus
        n += pas
        return n
    return inc            # inc garde n en vie
c = compteur(); c(); c()  # 1, 2
fs = [lambda: i for i in range(3)]
[f() for f in fs]         # [2, 2, 2] : liaison
fs = [lambda i=i: i for i in range(3)]  # [0, 1, 2]
\end{code}
\note{Une fonction est un objet : on la passe, la rend, la range dans un dict. Une variable libre est lue à l'appel, pas à la création.}
\end{bloc}

\begin{bloc}{Signatures}{syn}
\begin{code}
def f(a, b, /, c, *, d): ...  # a, b : par position seule
f(1, 2, 3, d=4); f(1, 2, c=3, d=4)  # d : par nom seul
\end{code}
\end{bloc}

\begin{bloc}{Décorateurs}{syn}
\begin{code}
import functools, time
def timer(f):            # f -> nouvelle f
    @functools.wraps(f)  # garde nom, doc
    def wrapper(*args, **kwargs):
        t0 = time.perf_counter()
        res = f(*args, **kwargs)
        dt = time.perf_counter() - t0
        print(f"{f.__name__} : {dt:.3f} s")
        return res
    return wrapper
@timer                   # run = timer(run)
def run(n): ...
def retry(n=3):          # avec paramètres :
    def deco(f):         # un niveau de plus
        @functools.wraps(f)
        def wrapper(*a, **kw):
            for k in range(n):
                try: return f(*a, **kw)
                except OSError:
                    if k == n - 1: raise
        return wrapper
    return deco
@retry(n=5)
def telecharger(url): ...
\end{code}
\note{Empilés, ils s'appliquent de bas en haut : \V|@a @b def f| donne \V|a(b(f))|.}
\end{bloc}

\begin{bloc}{functools}{lib}
\begin{code}
from functools import cache, lru_cache, partial, reduce
@cache                           # mémoïsation
def fib(n): return n if n < 2 else fib(n-1) + fib(n-2)
@lru_cache(maxsize=256)          # les 256 derniers
def cout(cfg): ...               # args hachables
hex8 = partial("{:08x}".format)  # hex8(42) : 0000002a
\end{code}
\note{Aussi : \V|@cached_property|, \V|@total_ordering|, \V|@singledispatch| (surcharge par type).}
\end{bloc}
\end{colonne}\colbreak
% ======================= colonne 3 : objets =======================
\begin{colonne}[objets][col-c]
\begin{bloc}{Méthodes spéciales}{typ}
\begin{optab}[27mm]
\V|__repr__ __str__|   & \V|repr(o)| pour déboguer, \V|str(o)| et \V|print| \\
\V|__eq__ __hash__|    & \V|==| ; redéfinir \V|__eq__| retire le hash \\
\V|__lt__ __le__...|   & \V|<|, \V|sorted|, \V|max| \\
\V|__len__ __bool__|   & \V|len(o)|, \V|if o:| \\
\V|__getitem__|        & \V|o[i]|, tranches ; rend itérable \\
\V|__iter__|           & \V|for x in o| ; \V|__contains__| : \V|x in o| \\
\V|__call__|           & \V|o(x)| : objet appelable \\
\V|__add__ __radd__|   & \V|o + x|, \V|x + o| ; \V|__iadd__| : \V|+=| \\
\V|__enter__ __exit__| & \V|with o:| (Context managers, au verso) \\
\V|__getattr__|        & attribut introuvable
\end{optab}
\end{bloc}

\begin{bloc}{Classes : compléments}{typ}
\begin{code}
class Temp:
    __slots__ = ("_c",)  # pas de __dict__
    def __init__(self, c): self.c = c
    @property
    def c(self): return self._c
    @c.setter            # t.c = 20 valide
    def c(self, v):
        if v < -273.15: raise ValueError(v)
        self._c = v
    @classmethod         # autre constructeur
    def depuis_f(cls, f): return cls((f - 32) / 1.8)
from abc import ABC, abstractmethod
class Bus(ABC):          # classe abstraite
    @abstractmethod
    def lire(self, adr: int) -> int: ...
\end{code}
\end{bloc}

\begin{bloc}{Héritage multiple}{typ}
\begin{code}
class Trace:  # mixin : un service
    def log(self, m): print(type(self).__name__, m)
class Uart(Trace, Periph): ...
Uart.__mro__  # ordre de recherche
\end{code}
\note{\V|super()| appelle la classe suivante du MRO, pas forcément la mère écrite : l'appeler dans chaque \V|__init__|.}
\end{bloc}

\begin{bloc}{dataclasses et Enum}{typ}
\begin{code}
from dataclasses import dataclass, field, asdict, replace
@dataclass(frozen=True, order=True, slots=True)
class Run:
    freq_mhz: float
    nom: str = "base"
    tags: list[str] = field(default_factory=list)
    def __post_init__(self):  # validation
        if self.freq_mhz <= 0: raise ValueError
r = Run(100.0); asdict(r); replace(r, nom="opt")
from enum import Enum, IntFlag, auto
class Etat(Enum):
    IDLE = auto(); RUN = auto()
Etat.RUN.name, Etat["IDLE"], Etat(1)
class Perm(IntFlag):
    R = 4; W = 2; X = 1
Perm.R | Perm.W               # combinable
\end{code}
\end{bloc}
\end{colonne}

%%%%%%%%%%%%%%%%%%%%%%%%%%%%%%%%%%%%%%%%%%%%%%%%%%%%%%%%%%%%%%%%%%%%%%%%%%%%%%
%                VERSO : TYPAGE, CONCURRENCE, QUALITÉ
%%%%%%%%%%%%%%%%%%%%%%%%%%%%%%%%%%%%%%%%%%%%%%%%%%%%%%%%%%%%%%%%%%%%%%%%%%%%%%
\newpage
\entete{Python avancé}{verso · typage, concurrence, qualité}
\noindent
% ======================= colonne 4 : with, typage, paquet =======================
\begin{colonne}[with · typage · paquet][col-a]
\begin{bloc}{Context managers}{ctl}
\begin{code}
class Chrono:
    def __enter__(self):
        self.t0 = time.perf_counter(); return self
    def __exit__(self, typ, val, tb):
        self.dt = time.perf_counter() - self.t0
        return False     # True : avale l'exc.
with Chrono() as c: travail()
print(c.dt)
from contextlib import contextmanager, suppress
@contextmanager
def dans(dossier):       # cd temporaire
    avant = os.getcwd(); os.chdir(dossier)
    try:
        yield            # corps du with
    finally:
        os.chdir(avant)  # même si exception
with dans("sim"), open("out.log", "w") as fh: ...
with suppress(FileNotFoundError): os.remove("x")
\end{code}
\note{\V|ExitStack| : nombre variable de \V|with| ; \V|nullcontext()| : \V|with| qui ne fait rien.}
\end{bloc}

\begin{bloc}{Annotations de type}{typ}
\begin{code}
from typing import Callable, Iterator, Literal, Protocol
def moy(xs: list[float]) -> float: ...
def trouver(k: str) -> int | None: ...  # Optional
cfg: dict[str, tuple[int, ...]] = {}
Mode = Literal["rtl", "gate"]
def appliquer(f: Callable[[int], str], x: int): ...
def premier[T](xs: list[T]) -> T:       # générique
    return xs[0]                        # 3.12
type Vecteur = list[float]              # alias 3.12
class Lisible(Protocol):                # duck typing
    def read(self, n: int = -1) -> bytes: ...
def charger(src: Lisible): ...          # a un .read()
def gen() -> Iterator[int]: yield 1     # génératrice
\end{code}
\note{Rien n'est vérifié à l'exécution : lancer \V|mypy «dossier»| ou \V|pyright|. Avant 3.12 : \V|T = TypeVar("T")|.}
\end{bloc}

\begin{bloc}{Autres structures typées}{typ}
\begin{code}
from typing import NamedTuple, TypedDict
class Pt(NamedTuple):  # tuple à champs nommés
    x: float; y: float
class Cfg(TypedDict):  # forme d'un dict (JSON)
    freq: float; seed: int
\end{code}
\end{bloc}

\begin{bloc}{Paquet installable}{lib}
\begin{code}
# pyproject.toml
[project]
name = "outils-sei"
version = "0.1.0"
dependencies = ["numpy"]
[project.scripts]
bilan = "outils_sei.bilan:main"  # commande bilan
\end{code}
\note{\V|pip install -e .| : installé en mode édition, \V|bilan| devient une commande.}
\end{bloc}
\end{colonne}\colbreak
% ======================= colonne 5 : concurrence =======================
\begin{colonne}[concurrence][col-b]
\begin{bloc}{Threads ou processus ?}{ex}
\begin{schema}[x=1mm, y=1mm, >={Stealth[length=1.5mm]}, line width=0.6pt,
    q/.style={draw=trait, rounded corners=1.2mm, fill=paper, inner sep=1.2mm, font=\fontbloc,
              align=center},
    r/.style={draw=cat, line width=0.9pt, fill=cat!12, inner sep=1.2mm, font=\fontbloc\bfseries,
              align=center, text width=31mm},
    o/.style={font=\fontbloc, text=note, inner sep=0.8pt}]
  \node[q] (q1) at (0,0)   {attend des E/S ?};
  \node[q] (q2) at (0,-8.5) {vectorisable ?};
  \node[r, anchor=west] (r1) at (22,0)    {threads ou asyncio};
  \node[r, anchor=west] (r2) at (22,-8.5) {numpy (libère le GIL)};
  \node[r, anchor=west] (r3) at (22,-17)  {processus : 1 par cœur};
  \draw[->] (q1) -- node[o, above] {oui} (r1);
  \draw[->] (q1) -- node[o, left] {non} (q2);
  \draw[->] (q2) -- node[o, above] {oui} (r2);
  \draw[->] (q2) |- node[o, left, pos=0.25] {non} (r3);
\end{schema}
\note{GIL : un seul thread exécute du Python à la fois. Sans GIL : Python 3.13 \V|python3.13t| (expérimental).}
\end{bloc}

\begin{bloc}{concurrent.futures : N simulations}{ex}
\begin{code}
import subprocess
from concurrent.futures import (ThreadPoolExecutor,
                                as_completed)
def simuler(seed: int) -> tuple[int, int]:
    r = subprocess.run(["vsim", "-c", "-do", "run.do",
                        "-sv_seed", str(seed)],
                       capture_output=True, text=True)
    return seed, r.returncode
with ThreadPoolExecutor(max_workers=8) as ex:
    futs = [ex.submit(simuler, s) for s in range(32)]
    for fut in as_completed(futs):  # ordre de fin
        seed, code = fut.result()
        print(seed, "OK" if code == 0 else "ÉCHEC")
\end{code}
\note{Threads : ils attendent des processus externes. \V|ex.map(f, args)| : résultats dans l'ordre.}
\end{bloc}

\begin{bloc}{Calcul en parallèle}{ex}
\begin{code}
from concurrent.futures import ProcessPoolExecutor
def cout(params): ...                  # pur calcul
if __name__ == "__main__":             # obligatoire
    with ProcessPoolExecutor() as ex:  # 1 par cœur
        res = list(ex.map(cout, grille, chunksize=16))
\end{code}
\note{Fonctions et arguments sont copiés (pickle) vers chaque processus : pas de lambda, pas d'état partagé.}
\end{bloc}

\begin{bloc}{asyncio}{ex}
\begin{code}
import asyncio
async def simuler(seed, sem):
    async with sem:                  # 8 à la fois max
        p = await asyncio.create_subprocess_exec(
            "vsim", "-c", "-do", "run.do",
            stdout=asyncio.subprocess.PIPE)
        out, _ = await p.communicate()
        return seed, p.returncode
async def main():
    sem = asyncio.Semaphore(8)
    res = await asyncio.gather(
        *(simuler(s, sem) for s in range(32)))
    async with asyncio.timeout(60):  # 3.11
        await asyncio.sleep(1)       # pas time.sleep
asyncio.run(main())                  # point d'entrée
\end{code}
\note{\V|async def| crée une coroutine ; \V|await| rend la main à la boucle d'événements. Un seul thread : un appel bloquant fige tout (\V|asyncio.to_thread(f)| pour l'isoler).}
\end{bloc}
\end{colonne}\colbreak
% ======================= colonne 6 : qualité =======================
\begin{colonne}[qualité][col-c]
\begin{bloc}{Tests avec pytest}{io}
\begin{code}
# tests/test_fifo.py : fichiers test_*.py
import pytest
from fifo import Fifo
def test_vide():
    assert Fifo(4).vide()     # assert simple
@pytest.fixture
def f():                      # préparation
    return Fifo(4)            # partagée
@pytest.mark.parametrize("n, plein",
                         [(3, False), (4, True)])
def test_plein(f, n, plein):  # f injectée
    for i in range(n): f.push(i)
    assert f.plein() == plein
def test_erreur(f):
    with pytest.raises(IndexError):
        f.pop()
assert 0.1 + 0.2 == pytest.approx(0.3)
\end{code}
\note{\V|pytest -x -v -k plein| : arrêt au 1\textsuperscript{er} échec, détaillé, filtre par nom. \V|tmp_path| : dossier temporaire fourni.}
\end{bloc}

\begin{bloc}{Débogage}{io}
\begin{code}
breakpoint()              # pause ici (pdb)
python3 -m pdb script.py  # dès la 1re ligne
import pdb; pdb.pm()      # après un plantage
import traceback; traceback.print_exc()
\end{code}
\begin{optab}[13mm]
\V|n s c| & ligne suivante, entrer dans l'appel, continuer \\
\V|p «x»| & afficher une expression ; \V|pp| joli \\
\V|l w q| & code autour, pile d'appels, quitter
\end{optab}
\end{bloc}

\begin{bloc}{Journal et mesures}{io}
\begin{code}
import logging
log = logging.getLogger(__name__)  # par module
logging.basicConfig(level=logging.INFO,
    format="%(asctime)s %(levelname)s %(message)s")
log.info("run %d", n); log.exception("échec")
python3 -m timeit "sum(range(10**6))"
python3 -m cProfile -s cumtime script.py
\end{code}
\note{\V|log.exception| ajoute la trace : à appeler dans un \V|except|.}
\end{bloc}

\begin{bloc}{Pièges avancés}{ex}
\begin{optab}[20mm]
Générateur épuisé & un 2\textsuperscript{e} parcours est vide : refaire l'appel ou \V|list(g)|. \\
Attribut de classe & \V|tags = []| au niveau classe : une seule liste pour toutes les instances. \\
\V|except:| nu & attrape aussi \V|KeyboardInterrupt| : viser des exceptions précises. \\
\V|__eq__| seul & l'objet n'est plus hachable (\V|set|, clé de \V|dict|). \\
Threads et CPU & aucun gain à cause du GIL : processus. \\
\V|await| oublié & \V|f()| seul crée la coroutine sans l'exécuter (« never awaited »).
\end{optab}
\end{bloc}
\end{colonne}
\end{document}
